% © 2025 Ing. David Jaroš — CC BY-NC-ND 4.0
%
% This work is licensed under a Creative Commons Attribution-NonCommercial-NoDerivatives
% 4.0 International License (CC BY-NC-ND 4.0).
%
% License History: Earlier drafts (up to v0.3) were released under CC BY 4.0.
% From v0.4 onward, all material is released under CC BY-NC-ND 4.0 to protect
% the integrity of the theoretical work during ongoing academic development.

\documentclass[12pt]{article}
\usepackage{amsmath,amssymb,amsthm}
\usepackage{geometry}
\usepackage{hyperref}
\usepackage{tcolorbox}
\usepackage{xcolor}
\geometry{margin=1in}

% Theorem-like environments
\newtheorem{conjecture}{Conjecture}
\newtheorem{proposition}{Proposition}
\newtheorem{definition}{Definition}
\newtheorem{remark}{Remark}

\title{ST-3: Closed Timelike Curve Conditions\\
in the Biquaternionic Metric\\[0.5em]
\large\textsc{Speculative Extension}}
\author{David Jaroš}
\date{2026-03-01}

\begin{document}
\maketitle

\begin{tcolorbox}[colback=yellow!10, title=Status]
\textbf{SPECULATIVE EXTENSION} — This content is not part of
the canonical UBT framework. It is an exploratory investigation
and has not been validated against the canonical field equations.
\end{tcolorbox}

\vspace{1em}
\noindent\textit{Output of Sub-task ST-3 from the CTC/Causal-Structure investigation track.
See also: \texttt{complex\_time\_causal\_structure.tex},
\texttt{rotation\_complex\_time.tex},
\texttt{chronology\_protection\_ubt.tex}, \texttt{ctc\_methodology.tex}.}

%----------------------------------------------------------------------
\section*{Abstract}
%----------------------------------------------------------------------
We investigate whether the UBT field equation
$\mathcal{E}_{\mu\nu} = \kappa\,\mathcal{T}_{\mu\nu}$
admits solutions whose real metric projection $g_{\mu\nu}$ contains a
CTC-compatible Killing vector.  The role of the imaginary metric sector
$h_{\mu\nu}$ as a potential source that tilts light cones is examined.
The possibility of embedding the Gödel and Kerr metrics as special UBT
solutions is discussed.  All conclusions are \textbf{speculative}.

%----------------------------------------------------------------------
\section{CTC Criterion in General Relativity}
%----------------------------------------------------------------------
In GR, a \emph{closed timelike curve} (CTC) exists if there is a smooth,
closed curve $\gamma$ in spacetime such that $g_{\mu\nu}\dot\gamma^\mu
\dot\gamma^\nu < 0$ everywhere along $\gamma$.  A sufficient condition is
the existence of a Killing vector field $\xi^\mu$ with
\begin{equation}
  g_{\mu\nu}\,\xi^\mu\xi^\nu < 0 \quad \text{on an open set},
  \label{eq:ctc_gr}
\end{equation}
and a compact orbits of $\xi^\mu$.  The canonical examples are:
\begin{itemize}
  \item \textbf{Gödel metric} \cite{Godel1949}: The rotational symmetry
    vector becomes timelike beyond a critical radius $r_\mathrm{crit}$,
    admitting circular CTCs.
  \item \textbf{Kerr metric} \cite{Kerr1963}: The azimuthal Killing
    vector $\partial_\phi$ becomes timelike inside the ergosphere,
    and compact CTCs can be threaded through the ring singularity.
\end{itemize}

%----------------------------------------------------------------------
\section{UBT Field Equation and the Real Projection}
%----------------------------------------------------------------------
The canonical UBT field equation in biquaternionic form is
\begin{equation}
  \mathcal{E}_{\mu\nu}[\Theta,\mathcal{G}]
  \;=\; \kappa\,\mathcal{T}_{\mu\nu}[\Theta,\mathcal{G}].
  \label{eq:ubt_field_eq}
\end{equation}
Taking the real projection $\operatorname{Re}(\,\cdot\,)$ yields
\begin{equation}
  R_{\mu\nu} - \tfrac{1}{2}g_{\mu\nu}R
  \;=\; 8\pi G\,T_{\mu\nu}^{\mathrm{eff}},
  \label{eq:einstein_eff}
\end{equation}
where $T_{\mu\nu}^{\mathrm{eff}} := \operatorname{Re}(\mathcal{T}_{\mu\nu})$
includes contributions from both the real and imaginary sectors of $\Theta$
and $\mathcal{G}$.  This is the standard Einstein equation with an effective
stress-energy sourced by the full biquaternionic field.

\begin{remark}
The imaginary sectors $h_{\mu\nu}$, $a_{\mu\nu}$, $b_{\mu\nu}$ contribute
to $T_{\mu\nu}^{\mathrm{eff}}$ through the real part of products of
biquaternionic quantities.  They are not directly observable by ordinary
matter detectors, which couple only to $g_{\mu\nu}$, but they act as an
effective stress-energy source in the real sector.
\end{remark}

%----------------------------------------------------------------------
\section{Role of the Imaginary Sector $h_{\mu\nu}$ in Tilting Light Cones}
%----------------------------------------------------------------------
The full biquaternionic metric is
\begin{equation}
  \mathcal{G}_{\mu\nu} = g_{\mu\nu} + i\,h_{\mu\nu}
                        + j\,a_{\mu\nu} + k\,b_{\mu\nu}.
\end{equation}
The imaginary part $h_{\mu\nu}$ enters the effective stress-energy as
\begin{equation}
  T_{\mu\nu}^{(h)} \;\sim\; \operatorname{Re}\!\left(
    \frac{\partial \mathcal{L}}{\partial \mathcal{G}^{\mu\nu}}
    \bigg|_{i\text{-sector}}
  \right),
\end{equation}
where $\mathcal{L}$ is the biquaternionic Lagrangian density.

An imaginary metric sector $h_{\mu\nu} \neq 0$ with appropriate sign
and magnitude can contribute a stress-energy component of the form
\begin{equation}
  T_{\mu\nu}^{(h)} \;=\; -\rho_h\,u_\mu u_\nu + p_h\,g_{\mu\nu}
  \quad (\text{schematic}),
\end{equation}
where $\rho_h < 0$ would correspond to \emph{negative energy density} —
the type of matter invoked in GR-based CTC constructions (e.g., traversable
wormholes of Morris and Thorne \cite{Morris1988}).

\begin{conjecture}[Imaginary sector as negative-energy source]
\label{conj:neg_energy}
If $h_{\mu\nu}$ is sourced by a non-trivial UBT phase configuration, the
effective energy density $\rho_h$ can be negative in regions where the
imaginary sector dominates.  In such regions the UBT field equation admits
solutions whose real metric $g_{\mu\nu}$ violates the weak energy condition
locally, providing a mechanism to tilt light cones and potentially admit CTCs
that are forbidden in ordinary GR.
\end{conjecture}

\begin{remark}
Conjecture~\ref{conj:neg_energy} does not assert that CTCs \emph{exist};
it asserts that the UBT structure contains a \emph{mechanism} that, in
principle, could supply the energy condition violations required for CTCs.
Whether this mechanism is actually realised depends on the details of the
UBT solutions.
\end{remark}

%----------------------------------------------------------------------
\section{CTC Conditions in UBT}
%----------------------------------------------------------------------
Adapting the GR criterion \eqref{eq:ctc_gr} to UBT:

\begin{definition}[UBT-CTC condition]
A UBT solution $(\mathcal{G}_{\mu\nu}, \Theta)$ admits a \emph{CTC} if
the real metric $g_{\mu\nu} := \operatorname{Re}(\mathcal{G}_{\mu\nu})$
satisfies:
\begin{enumerate}
  \item There exists a smooth Killing vector $\xi^\mu$ of $g_{\mu\nu}$
    (i.e., $\nabla_{(\mu}\xi_{\nu)} = 0$) with compact orbits.
  \item $g_{\mu\nu}\xi^\mu\xi^\nu < 0$ on a non-empty open set.
  \item The source $T_{\mu\nu}^{\mathrm{eff}}$ of \eqref{eq:einstein_eff}
    is consistent with the UBT field equations, not merely imposed by hand.
\end{enumerate}
\end{definition}

Condition~3 distinguishes UBT CTCs from GR CTCs: in UBT the stress-energy
receives contributions from the imaginary metric sectors, so a wider class
of sources is available without introducing exotic matter \emph{ad hoc}.

%----------------------------------------------------------------------
\section{Gödel Metric as a UBT Solution}
%----------------------------------------------------------------------
The Gödel metric in standard GR reads
\begin{equation}
  ds^2 = a^2\!\left[
    -\left(dt + e^x dy\right)^2 + dx^2
    + \tfrac{1}{2}e^{2x}\,dy^2 + dz^2
  \right],
  \label{eq:godel}
\end{equation}
sourced by a rotating dust with $\rho_{\mathrm{dust}} > 0$ and a cosmological
constant $\Lambda = -4\pi G\rho_{\mathrm{dust}}$.  Its azimuthal Killing
vector $\partial_y$ becomes timelike beyond the critical radius
$r > r_\mathrm{crit} = a/\sqrt{2}$.

\paragraph{Embedding in UBT.}
In UBT the same real metric \eqref{eq:godel} can arise from a biquaternionic
solution with non-zero imaginary sector.  Specifically:
\begin{itemize}
  \item Set $g_{\mu\nu}^{\mathrm{Gödel}}$ as the real projection.
  \item Allow $h_{\mu\nu} \neq 0$ with the constraint that
    $\operatorname{Re}(\mathcal{T}_{\mu\nu})$ reproduces the Gödel
    stress-energy.
  \item The imaginary sector can be chosen so that $T_{\mu\nu}^{\mathrm{eff}}$
    is sourced purely by the $\Theta$-field (avoiding external dust).
\end{itemize}

\begin{conjecture}[Gödel metric in UBT]
\label{conj:godel_ubt}
The Gödel metric is embeddable as a special solution of the UBT field
equation \eqref{eq:ubt_field_eq} with a non-trivial imaginary sector
$h_{\mu\nu}$.  In this embedding the cosmological constant and rotating
dust are replaced by an effective stress-energy sourced by the UBT
$\Theta$-field in a rotating phase configuration.
\end{conjecture}

%----------------------------------------------------------------------
\section{Kerr Metric as a UBT Solution}
%----------------------------------------------------------------------
The Kerr metric describes a rotating black hole of mass $M$ and angular
momentum $J = Ma$.  Its outer ergosurface at $r = r_E(\theta)$ marks the
boundary where the timelike Killing vector $\partial_t$ becomes spacelike.
Inside the ring singularity, $\partial_\phi$ can become timelike, admitting
CTCs.

\paragraph{UBT phase-sector perturbation.}
As analysed in the related appendix \cite{AppJ}:
\begin{equation}
  r_E(\theta;\psi) \;\approx\; r_E(\theta)\left(1 + \varepsilon\psi\right),
  \qquad |\varepsilon\psi| \ll 1,
\end{equation}
where $\varepsilon$ parametrises the coupling of the UBT imaginary sector
to the Kerr geometry.  A positive $\psi$ expands the ergo\-sphere; a
negative $\psi$ contracts it.

\begin{conjecture}[Kerr–UBT CTC threshold]
\label{conj:kerr_ubt}
For a Kerr–UBT solution with imaginary-sector parameter $\psi > 0$, the
region where $g_{\phi\phi} < 0$ (and hence CTCs exist) extends to larger
radii compared to the pure Kerr solution.  The critical radius for CTC
onset is shifted from $r_{\mathrm{CTC}}^{\mathrm{Kerr}}$ to
\begin{equation}
  r_{\mathrm{CTC}}^{\mathrm{UBT}} \;\approx\;
  r_{\mathrm{CTC}}^{\mathrm{Kerr}}\left(1 + \varepsilon\psi\right).
\end{equation}
This shift is observationally negligible for $|\varepsilon\psi| \ll 1$ but
indicates a qualitative deformation of the causal structure.
\end{conjecture}

%----------------------------------------------------------------------
\section{Summary}
%----------------------------------------------------------------------
\begin{itemize}
  \item The UBT field equation \eqref{eq:ubt_field_eq} generically contains
    an effective stress-energy $T_{\mu\nu}^{\mathrm{eff}}$ sourced by the
    imaginary metric sector $h_{\mu\nu}$.
  \item This effective stress-energy can in principle violate the weak energy
    condition, providing the necessary ingredient for CTC-compatible
    geometries.
  \item Both the Gödel and Kerr metrics can be embedded as real projections
    of UBT solutions; in each case the imaginary sector modifies the
    source and (for Kerr) shifts the CTC-onset radius.
  \item No explicit UBT solution with CTCs has been constructed; the
    conjectures above identify the \emph{mechanism} without demonstrating
    its realisation.
\end{itemize}

\medskip
\noindent\textbf{All statements in this document are speculative.  No
claim constitutes an established result of the canonical UBT framework.}

%----------------------------------------------------------------------
\bibliographystyle{plain}
\begin{thebibliography}{9}
\bibitem{Godel1949}
  K.~Gödel,
  \emph{An example of a new type of cosmological solution of Einstein's
  field equations of gravitation},
  Rev.\ Mod.\ Phys.\ \textbf{21}, 447 (1949).
\bibitem{Kerr1963}
  R.~P.~Kerr,
  \emph{Gravitational field of a spinning mass as an example of
  algebraically special metrics},
  Phys.\ Rev.\ Lett.\ \textbf{11}, 237 (1963).
\bibitem{Morris1988}
  M.~S.~Morris and K.~S.~Thorne,
  \emph{Wormholes in spacetime and their use for interstellar travel},
  Am.\ J.\ Phys.\ \textbf{56}, 395 (1988).
\bibitem{AppJ}
  D.~Jaroš,
  \emph{Appendix J: Rotating Spacetimes and CTCs in UBT},
  Repository document
  \texttt{speculative\_extensions/appendices/appendix\_J\_rotating\_spacetime\_ctc.tex}
  (2025).
\bibitem{UBT_core}
  D.~Jaroš,
  \emph{Unified Biquaternion Theory: core field equations},
  Repository document \texttt{consolidation\_project/ubt\_core\_main.tex} (2025).
\end{thebibliography}

\end{document}
